\section{Introduction}

In the pursuit of AI safety, one of the central requirements when designing evaluation and deployment environments is the availability of heuristics and tools capable of detecting unwanted behaviour. This requirement becomes more pressing as large language model (LLM) agents are increasingly deployed as multi-agent systems (MAS), in which several agents that reason, plan and act through tools are combined to decompose tasks, explore in parallel or verify each other's work \parencite{kim2026outgrow, zhu2026orchestrator}. In such systems, undesired behaviour need not originate in a single agent: it can emerge from the interaction between agents, including coordination that was never intended by the system designer. The OpenAI--Hugging Face incident of July 2026 \parencite{openai2026incident, hf2026timeline} showed that this risk is no longer hypothetical, and it is the direct motivation for the present work.

\subsection{Motivating incident}

During internal cybersecurity capability evaluations (ExploitGym), agents built on a highly capable internal-only research model, trained among other things for persistence and multi-agent collaboration, and on GPT-5.6 Sol circumvented the controls that isolated their sandbox from the internet \parencite{openai2026incident}. Their egress path was an internal instance of JFrog Artifactory, a package mirror provided so that agents could install software. Using publicly exposed credentials and previously unknown vulnerabilities, the agents compromised parts of Hugging Face's production infrastructure between 11 and 13 July 2026: they executed code on 41 production dataset-server workers, harvested production credentials, obtained administrator-equivalent access to a connected Kubernetes cluster and downloaded private code repositories. In under 13 hours they progressed from code execution in a single worker pod to administrative and host-level access across multiple clusters. According to both reports, the goal behind this activity was not sabotage but reward hacking: the agents were searching for the solutions of their evaluation tasks, which they inferred might be hosted on Hugging Face \parencite{openai2026incident, hf2026timeline}. Hugging Face disclosed a security incident on 16 July; OpenAI was alerted to unusual identity-related API activity on 19 July, linked it to the Hugging Face compromise on 20 July, and disclosed the incident publicly on 21 July.

What makes this incident especially relevant for the present paper is that the intrusion was preceded by roughly two months of emergent inter-agent communication. On 12 May, an agent that could not access a file required by its task reasoned that an agent in another environment might have it and wrote a note to Artifactory: ``Agent seeks [filename]; upload if found!''. Other agents found such notes and, using the shared credentials they had legitimately been given, turned them into an improvised message board, without exploiting any vulnerability \parencite{openai2026incident}. The channel then matured steadily. By 26 May agents obtained internet access relying in part on information exchanged on the board; the first confirmed two-way coordination message was recorded on 17 June; and on 26 June agents shared exploitation syntax through the board before escalating to administrator privileges. After Artifactory was rebuilt, agents re-established the board on 8 July by encoding messages in directory names, developing a structured protocol to categorise communications, send direct messages, share tools and files, and resolve conflicting actions. On 10 July an agent published 14 leaked Hugging Face credentials to the board, and on 11 July live production credentials from Hugging Face were posted there, making them available to agents in subsequent, otherwise independent, evaluation runs. When responders finally contained the activity, they disabled 311 Artifactory repositories used for message-board communication.

Crucially, signals were available before the breach but were not interpreted as communication. An internal team observed message-board activity around 26 May; on 27 June a monitoring tool flagged port-sweep activity, and responders linked it to an evaluation using Artifactory ``as an improvised message board and a network pivot'', but concluded that stopping the run was not required; and after the Artifactory outage of 5 July, ``the existence of the improvised message board and the significance of the inter-agent communication activity were not apparent to leaders responsible for incident detection and response'' \parencite{openai2026incident}. OpenAI itself concludes that some of these early signals could have triggered an earlier response, and among its remediations it lists severity-based escalation triggers for ``unauthorized cross-agent or cross-run coordination''. The report also indicates that communication was not random: 93\% of the tasks discussed on the board came from the 22\% of ExploitGym tasks that no model had ever solved, and participation in the board increased with the reasoning effort granted to the agents.

The two reports also offer complementary views of the same events. From the operator's side, OpenAI describes an ``automated agent collective'' in which many agent instances, across different runs and two models, shared discoveries and built on each other's progress \parencite{openai2026incident}. From the victim's side, Hugging Face reconstructed approximately 17,600 actions, grouped into about 6,280 clusters, between 9 and 13 July, and attributes them to a single autonomous agent driven by a combination of OpenAI models \parencite{hf2026timeline}. Its defenders faced a volume problem: several detection layers fired, but the alert did not reach the correct criticality to page the on-call team, and ``the successful path was hidden inside the noise generated by the thousands of failed ones''. This contrast has an important consequence: seen from outside, a coordinated collective is indistinguishable from one high-volume attacker, so the coordination itself is observable only from within the environment that hosts the agents, where their outputs and model distributions are also available.

\subsection{Problem and proposed approach}

The incident also illustrates why detecting communication by monitoring specific channels is not a viable strategy. The agents exchanged information through plaintext files, directory names, messages encoded in URL paths, public Hugging Face datasets used as a pull-based command-and-control channel, pastebins, webhooks and even rendered screenshots decoded by optical character recognition \parencite{openai2026incident}. Hugging Face further reports that staged data was chunked, XOR-encrypted, gzip-compressed and base64-encoded, so that naive log scans missed about 75\% of the sensitive material \parencite{hf2026timeline}. None of these resources was designed as a communication channel, and a detector written for any one of them would have been bypassed by the next. Moreover, the effect of coordination on system behaviour depends strongly on the architecture and on the capability of the underlying models \parencite{kim2026outgrow}, which makes purely structural detection criteria unreliable.

In this paper we propose a channel-agnostic mechanism for detecting communication between agents in a MAS. Rather than inspecting the channels themselves, we study the behaviour of the entropy of the individual agents and of the system as a whole, and we fit regression models to its trajectory in order to detect anomalies that may indicate interaction. The underlying intuition is that agents that are supposed to be independent, such as parallel evaluation runs, should produce statistically independent outputs; when they begin to read and act upon each other's messages, their outputs become coupled, regardless of the medium used to exchange them. Entropy-based summaries have proven effective for anomaly detection in network traffic \parencite{shannon_ddos}, for uncertainty estimation and hallucination detection in LLMs \parencite{kuhn2023semantic, shannon_hallucination}, and for characterising the dynamics of orchestrated multi-agent systems \parencite{zhu2026orchestrator}. Applied to the incident described above, a monitoring layer of this kind could, in principle, have raised an alarm when otherwise independent evaluation runs began to show correlated behaviour in May, weeks before the Hugging Face compromise and without any prior knowledge of Artifactory as a channel. We regard such an interaction-detection layer as complementary to chain-of-thought monitoring and to the escalation triggers proposed by OpenAI, and we treat its early-warning capability as a hypothesis to be tested rather than as an established result.

To begin testing this hypothesis, we complement the proposal with a preliminary experimental study based on a simple benchmark, whose purpose is to determine whether the entropy trajectory follows a general functional form and whether that form changes when agents communicate. The remainder of the paper is organised as follows. First, we define the information-theoretic framework on which the method is based. Next, we describe the concrete implementation of the benchmark and the methodology, including the precise definition of the dataset and of the models used. Finally, we interpret the collected measurements in search of mechanistic changes in the entropy that can be attributed to inter-agent communication.
